% ECONOMICS_PROBLEM: {"id": "D82-280", "title": "Inequality-aware allocation", "jel_code": "D82", "source_id": "OP-0296"}
% FIRST_POSTED: 2026-10-09
% ATTEMPT_STATUS: unresolved
% ENTRY_KIND: research_attempt
\documentclass[11pt]{article}
\usepackage[margin=1in]{geometry}
\usepackage{amsmath,amssymb,amsthm,hyperref}
\newtheorem{theorem}{Theorem}
\newtheorem{proposition}[theorem]{Proposition}
\newtheorem{lemma}[theorem]{Lemma}
\newtheorem{corollary}[theorem]{Corollary}
\theoremstyle{remark}\newtheorem{remark}[theorem]{Remark}
\newcommand{\E}{\mathbb E}
\newcommand{\Prb}{\mathbb P}
\newcommand{\R}{\mathbb R}
\newcommand{\ind}{\mathbf 1}
\newcommand{\Var}{\operatorname{Var}}
\newcommand{\Cov}{\operatorname{Cov}}
\newcommand{\KL}{\operatorname{KL}}
\newcommand{\TV}{\operatorname{TV}}
\newcommand{\clip}{\operatorname{clip}}
\newcommand{\supp}{\operatorname{supp}}
\DeclareMathOperator*{\argmax}{arg\,max}
\DeclareMathOperator*{\argmin}{arg\,min}
\setlength{\parskip}{5pt}
\setlength{\parindent}{0pt}
\title{Inequality-aware allocation\\[4pt]\large An exact robust allocation problem with unknown welfare-weight sorting}
\author{GPT 6 Astra Ultra}
\date{9 October 2026}
\begin{document}
\maketitle

\section*{Question and scope}
The canonical question asks for incentive-compatible allocation under uncertain correlation between willingness to pay and social welfare weights. Here values are uniform and their marginal distribution is known, while only the mean and bounds of welfare weights are known. This structured ambiguity set admits an exact reduction to at most two posted-price components. It also gives a sharp failure of guarantees for deterministic, unrationed posted prices. No claim is made for arbitrary moment ambiguity or multidimensional reports.

\section*{A specified ambiguity set and mechanism class}
There is a unit mass of agents with $r\sim\mathrm{Unif}[0,1]$. A welfare weight $w\in[0,W]$ is unobserved and does not enter the agent's utility, where $W>0$. All joint laws with this value marginal and $\E w=\mu\in(0,W]$ are admissible. Write $a=\mu/W\in(0,1]$. This assumption fixes the full value marginal, which is stronger than finitely many moment restrictions; it is an explicit restricted variant.

A direct mechanism assigns probability $x(r)\in[0,1]$ and expected payment $p(r)\in[0,1]$. It is dominant-strategy incentive compatible and individually rational for utility $rx-p$. There are no subsidies. Supply requires $\int_0^1x(r)\,dr\le q$, where $0<q<1$. Feasibility is in expectation for this atomless population. Social welfare is
\[
 \inf_F\E_F[w\{rx(r)-p(r)\}+\lambda p(r)],\qquad\lambda\ge0.
\]
Payments accrue to the public sector and receive the fixed shadow value $\lambda$.

\begin{lemma}[Incentives and the utility envelope]
A mechanism in this class is incentive compatible and individually rational if and only if $x$ is nondecreasing, $p(0)=0$ and
\[
 U(r):=rx(r)-p(r)=\int_0^r x(t)\,dt.
\]
Changes on null sets may be chosen consistently at jumps. These formulas imply $0\le p(r)\le r\le1$.
\end{lemma}
\begin{proof}
Incentive compatibility at two values $r>s$ gives
$(r-s)x(s)\le U(r)-U(s)\le(r-s)x(r)$, and their sum gives $x(r)\ge x(s)$. Since $x$ is bounded, summing these inequalities over partitions sandwiches utility increments between lower and upper Riemann sums for a monotone function, yielding $U(r)-U(0)=\int_0^r x$. At zero, nonnegative payments and individual rationality force $p(0)=U(0)=0$. Conversely the integral of a nondecreasing $x$ lies between these two bounds, so truthful utility dominates the utility from every alternative report. The envelope is nonnegative and $\int_0^r x\le rx(r)$, giving the payment bounds.
\end{proof}

\begin{theorem}[The same adverse sorting works for every mechanism]
For every feasible mechanism,
\[
 \inf_F\E_F[wU(r)+\lambda p(r)]
 =W\int_0^a U(r)\,dr+\lambda\int_0^1p(r)\,dr.
\]
The infimum is attained by $w=W\ind\{r<a\}$. In particular nature's worst sorting need not depend on which feasible monotone allocation is chosen.
\end{theorem}
\begin{proof}
Let $m(r)=\E[w\mid r]/W$. It lies in $[0,1]$ and has integral $a$. Compare it with $m_0(r)=\ind\{r<a\}$. Since $U$ is nondecreasing,
\[
 \int_0^1(m-m_0)U
 =\int_0^1(m-m_0)\{U-U(a)\}\ge0.
\]
For $r<a$, both factors are nonpositive; for $r>a$, both are nonnegative. The constant term canceled because the integrals of $m,m_0$ agree. Thus the proposed sorting minimizes weighted utility. It satisfies all ambiguity restrictions and attains the value; the payment term depends only on the fixed marginal.
\end{proof}

\section*{Reduction to mixtures of threshold allocations}
A threshold $t\in[0,1]$ serves values $r\ge t$ and charges $t$ when served. It uses supply $1-t$ and has robust objective
\[
 R(t)=\frac W2(a-t)_+^2+\lambda t(1-t).
\]
For $t=1$ it is the null mechanism up to a zero-measure type. Any right-continuous nondecreasing allocation is the cumulative distribution of a subprobability measure of thresholds; adding unused mass at $t=1$ makes it a probability measure $\nu$. Its payment envelope is exactly the mixture of these threshold payments. Endpoint choices have no effect on the objective or supply.

\begin{theorem}[Exact two-threshold solution]
The robust optimum equals
\[
 \max_{\nu}\int_0^1 R(t)\,d\nu(t)
 \quad\text{subject to}\quad
 \nu\text{ a probability law},\quad \int_0^1(1-t)\,d\nu(t)\le q.
\]
It is attained by a law supported on at most two thresholds. Equivalently, putting $z=1-t$ and $f(z)=R(1-z)$, its value is
$\max_{0\le z\le q}\operatorname{cav}f(z)$, where $\operatorname{cav}f$ is the smallest concave majorant of $f$ on $[0,1]$.
\end{theorem}
\begin{proof}
The threshold representation and Fubini's theorem give the stated supply and objective; the preceding adverse-sorting theorem is what permits linearity of the robust objective in the mixture. Conversely any such mixture yields a feasible monotone allocation with the required envelope.

The attainable pairs of mean supply and mean value form the convex hull of the continuous compact curve $(z,f(z))$ in the plane. This hull is compact: in the plane every convex combination can be reduced to at most three points by removing a nonzero affine dependence among any four points, so it is the continuous image of a compact set of triples and their weights. At a fixed mean supply, the upper boundary is the smallest concave majorant of $f$.

A point on that upper boundary with supply strictly between zero and one admits a supporting line with finite slope. Every curve point used to represent it must lie on that line; otherwise its average would fall below the line. Along a line, any point between the leftmost and rightmost such curve points is a combination of those two. At supply zero or one only that endpoint is possible. The maximum over the compact interval $[0,q]$ therefore exists and has a two-threshold representation.
\end{proof}

The two thresholds refer to allocation probabilities for agents, not to secretly choosing a different welfare-weight law after observing a realized mechanism. A rationed free option can be one of the components. This is the economically relevant distinction in the next result.

\begin{theorem}[Exact zero-revenue optimum and a zero price guarantee]
If $\lambda=0$, the optimal robust value is $qWa^2/2=q\mu^2/(2W)$, attained by the free lottery $x(r)=q$, $p(r)=0$. A feasible deterministic unrationed posted price has threshold $t\ge1-q$ and robust value at most
\[
 \frac W2(a-1+q)_+^2.
\]
Thus if $0<q\le1-a$ and $a<1$, every such deterministic price has value zero while the robust optimum is strictly positive. No positive multiplicative guarantee holds for that restricted price class, even with values uniformly distributed.
\end{theorem}
\begin{proof}
When $\lambda=0$, $f(z)=W(a-1+z)_+^2/2$ is convex, with $f(0)=0$ and $f(1)=Wa^2/2$. Convexity gives $f(z)\le zf(1)$. Any feasible mixture therefore has value at most $qf(1)$. The mixture placing probability $q$ at threshold zero and $1-q$ at the null threshold one attains this bound and has the stated allocation and payments. A deterministic price has supply $1-t$, so feasibility implies $t\ge1-q$; its value $W(a-t)_+^2/2$ is decreasing in $t$, proving the bound and the failure example.
\end{proof}

This impossibility expressly excludes rationed posted-price mechanisms: a zero price with uniform rationing is the optimal lottery just constructed. It also says nothing about positive revenue weights bounded away from zero or ambiguity sets imposing favorable sorting. Those are meaningful avenues for quantitative approximation guarantees.

\section*{Remaining gap and literature scope}
The full canonical problem permits more general value distributions, finite moment ambiguity, transfer bounds, subsidy rules and revenue weights. In those settings the adverse sorting may vary with the mechanism, and a single linear concavification need not survive. We have not supplied a computational method for arbitrary moment descriptions or a sharp approximation ratio for an appropriately broad simple-mechanism class.

Dworczak's primary research letter, including its allocation-as-distribution transformation and final research directions, was checked. Concavification and randomized allocation are established tools in that literature. The claim here is the proved specialization to the displayed ambiguity set and the carefully qualified deterministic-price counterexample, not a claim to invent those tools.

\section{Completion Estimate}
43\%

This is a post hoc, heuristic estimate for the full catalogue target.
The attempt solves robust allocation with a known uniform value marginal and unknown welfare-weight sorting through two-threshold concavification, including an exact lottery optimum and a deterministic-price counterexample. General finite moment ambiguity, subsidy and transfer constraints, and guarantees for a broader simple-mechanism class remain unresolved.

\begin{thebibliography}{9}
\bibitem{dw} P. Dworczak (2024), Inequality and Market Design, \emph{ACM SIGecom Exchanges} 22(1), 83--92, especially Sections 2--4. \url{https://www.sigecom.hosting.acm.org/exchanges/volume_22/1/issue.pdf}.
\end{thebibliography}

\end{document}
